\cotefile{79}
\scanpage{17}

où

\[
\Lambda_0=\underset{\scriptstyle n\longmapsto n.1_\Lambda}{\operatorname{Im}(\mathbb{Z}\longrightarrow\Lambda)}
\tag*{(39)}
\]

donc l'ens des composantes connexes est l'ens. \deleted{des} quotient

\[
L_x/\Lambda_0=L_x\mathbin{\wedge}_{\Lambda}\Lambda/\Lambda_0\;.\tag*{(40)}
\]

Ainsi, la structure $SP_M$ est une vraie structure \inserted{de $1$-polyèdre} triangul\deleted{\unsure{ée}}é, \deleted{\unsure{ssmt|satisfait}} i.e. satisfait à l'axiome de connexité pour le contour associé à un sommet, sss $\Lambda_0=\Lambda$, i.e. sss on a

\[
\Lambda\simeq\mathbb{Z}\quad\text{ou}\quad\unsure{\rightarrow}\ \Lambda\simeq\mathbb{Z}/n\mathbb{Z}\tag*{(41)}
\]

pour $n\in\mathbb{N}^{\deleted{\unsure{*}}}$, $n\geq 3$ (\uline{NB} on a exclu les cas $n=1$ ou $2$, avec la condition $2.1_\Lambda\ne 0$)

\uline{Remarque} Si on avait travaillé dans $P_M$ au lieu de $SP_M$, on trouverait que \deleted{\illegible{1}} l'ens des sommets liés à $s=\dot{x}$ \inserted{dans $P_M$} \deleted{est en corr. $1$-$1$ avec l'ens quotient} \inserted{admet l'ens \unsure{précédent|précident} comme} \unsure{comp.} \unsure{connexes|connexe} -- donc on a connexité sss on

% archive sheet number: 16
